\section{Berechenbarkeitstheorie}

\subsection{Church-Turing-These}
\textbf{Intuitiv berechenbar:} Funktionen, die durch ein mechanisches Verfahren berechnet werden können.\\
\textbf{Turing-berechenbar:} Funktionen, die von einer TM berechnet werden können.\\
\textbf{Church-Turing-These:} Die Klasse der Turing-berechenbaren Funktionen stimmt mit der Klasse der intuitiv berechenbaren Funktionen überein.\\
\textbf{Gandy's These M:} Alles, was mit einer endlichen Maschine berechnet werden kann, ist bereits Turing-berechenbar. 
\textit{Nicht formal beweisbar.}

\subsection{Turing-berechenbare Funktionen}
Eine TM $T = (Q, \Sigma, \Gamma, \ldots)$ als partielle Funktion $T : \Sigma^* \to \Gamma^*$:
\[
    T(w) = \begin{cases} u & \text{falls } T \text{ auf } w \text{ mit } u \text{ auf dem Band anhält} \\ \uparrow & \text{falls } T \text{ bei Input } w \text{ nicht anhält} \end{cases}
\]
$f : \Sigma^* \to \Gamma^*$ heisst \textbf{Turing-berechenbar}, wenn es eine TM $T$ gibt mit $T(w) = f(w)$ für alle $w$ wo $f$ definiert ist.\\

\subsection{LOOP-Programme}
\textbf{Syntax:}
\begin{itemize}
    \item \textbf{Zuweisungen:} \texttt{x = y + c} und \texttt{x = y - c} wobei $x, y$ Variablen und $c \in \mathbb{N}$ ist.
    \item \textbf{Sequenzen:} \texttt{P ; Q}
    \item \textbf{Schleifen:} \texttt{Loop x Do P End}
\end{itemize}
\textbf{Semantik:} $P_k(n_1, \ldots, n_k) =$ Wert von $x_0$ nach Ablauf, mit $n_i$ in $x_i$.\\
\textbf{Konventionen:} $x_0$ und Hilfsvariablen mit 0 initialisiert. Subtraktion $\geq 0$.

\subsubsection{Bedingte Ausführung (IF x1=0 THEN P)}
\begin{verbatim}
x2 = x2 + 1;
Loop x1 Do x2 = x2 - 1 End;
Loop x2 Do P End
\end{verbatim}

\subsubsection{Verzweigung (IF x1=0 THEN P ELSE Q)}
\begin{verbatim}
x2 = x2 + 1; x3 = x3 + 1;
Loop x1 Do x2 = x2 - 1; x4 = x3 + 0 End;
Loop x2 Do P End;
Loop x4 Do Q End
\end{verbatim}

\subsection{WHILE-Programme}
Erweitern LOOP um: \texttt{While x > 0 Do P End}\\
\begin{itemize}
    \item Jedes LOOP-Programm ist ein WHILE-Programm
    \item WHILE-Programme terminieren \textbf{nicht} immer
\end{itemize}

\subsection{GOTO-Programme}
\textbf{Syntax:}
\begin{itemize}
    \item \texttt{Mk: xi = xj + c} / \texttt{Mk: xi = xj - c}
    \item \texttt{Mk: If xi = c Then Goto Mr}
    \item \texttt{Mk: Goto Mr}
    \item \texttt{Mk: HALT}
    \item Sequenzen: \texttt{P ; Q} (keine gemeinsamen Marker)
\end{itemize}

\subsection{Primitiv rekursive Funktionen}
\subsubsection{Grundfunktionen}
\begin{itemize}
    \item \textbf{Konstante:} $c_k^n : \mathbb{N}^n \to \mathbb{N}$, $c_k^n(x_1, \ldots, x_n) = k$
    \item \textbf{Nachfolger:} $\eta : \mathbb{N} \to \mathbb{N}$, $\eta(x) = x + 1$
    \item \textbf{Projektion:} $\pi_k^n : \mathbb{N}^n \to \mathbb{N}$, $\pi_k^n(x_1, \ldots, x_n) = x_k$
\end{itemize}

\subsubsection{Abschluss durch Einsetzung}
Sind $f, g_1, \ldots, g_k$ primitiv rekursiv mit $g_i : \mathbb{N}^n \to \mathbb{N}$, $f : \mathbb{N}^k \to \mathbb{N}$, dann auch:
\[
    h(\vec{x}) = f(g_1(\vec{x}), \ldots, g_k(\vec{x}))
\]

\subsubsection{Abschluss durch primitive Rekursion}
Sind $h : \mathbb{N}^n \to \mathbb{N}$ und $g : \mathbb{N}^{n+2} \to \mathbb{N}$ primitiv rekursiv, dann auch $f : \mathbb{N}^{n+1} \to \mathbb{N}$ mit:
\begin{align*}
    f(0, \vec{x}) &= h(\vec{x}) \\
    f(k+1, \vec{x}) &= g(f(k, \vec{x}), k, \vec{x})
\end{align*}

\subsubsection{Beispiele}
\textbf{Addition:}
$\text{Add}(0,y) = \pi_1^1(y)$, \quad $\text{Add}(x+1,y) = \eta(\pi_1^3(\text{Add}(x,y),x,y))$\\
\textbf{Multiplikation:}
$\text{Mult}(0,y) = c_0^1(y)$, \quad $\text{Mult}(x+1,y) = \text{Add}(\text{Mult}(x,y), \pi_2^2(x,y))$

\subsection{Unvollständigkeit von LOOP / Prim. Rek.}
LOOP-Programme und primitiv rekursive Funktionen sind \textbf{nicht} Turing-vollständig. Es gibt totale, berechenbare Funktionen, die nicht LOOP-berechenbar sind:

\subsubsection{Ackermannfunktion}
\begin{align*}
    a(0, m) &= m + 1 \\
    a(n+1, 0) &= a(n, 1) \\
    a(n+1, m+1) &= a(n, a(n+1, m))
\end{align*}
\textbf{Eigenschaften:}
\begin{itemize}
    \item Turing-berechenbar (WHILE-Programm möglich)
    \item \textbf{Nicht} primitiv rekursiv / LOOP-berechenbar
    \item Total (überall definiert) -- Beweis via Induktion über lexikograph.\ Ordnung auf $\mathbb{N} \times \mathbb{N}$
    \item Wächst schneller als jede primitiv rekursive Funktion
\end{itemize}

\subsubsection{LOOP-Interpreter}
$I : \mathbb{N}^2 \to \mathbb{N}$ mit $I(\langle P \rangle, x) = P_1(x)$ für jedes LOOP-Programm $P$.\\
\textbf{Satz:}
\begin{itemize}
    \item $I$ ist total und eindeutig (LOOP terminiert immer)
    \item $I$ ist Turing-berechenbar
    \item $I$ ist \textbf{nicht} LOOP-berechenbar
\end{itemize}
\textbf{Beweis (Widerspruch):} Angenommen $P_2$ ist LOOP-Interpreter. Definiere LOOP-Programm $Q$ mit $Q_1(x) = P_2(x,x) + 1$. Dann:
\[
    Q_1(\langle Q \rangle) = P_2(\langle Q \rangle, \langle Q \rangle) + 1 = Q_1(\langle Q \rangle) + 1 \quad \lightning
\]

\subsection{Überblick Berechnungsmodelle}
\[
    \underbrace{\text{LOOP} = \text{Prim.\ rek.}}_{\text{nicht Turing-vollständig}} \subsetneq \underbrace{\text{TM} = \text{WHILE} = \text{GOTO}}_{\text{Turing-vollständig}}
\]
